\pg{438} b) Todos os homens, em todos os tempos e em todos os logares, reconheceram e admittiram a existencia de um Ente supremo. Este consenso unanime, constante e universal, não póde ser effeito do \emph{erro}; pois o \emph{erro} é causa de discordias, é mutavel e particular. Logo é effeito da \emph{verdade}; porque só a \emph{verdade} é unanime, constante e universal, e só ella pode levar os homens a mortificar as paixões e renunciar á propria vontade, para se sujeitarem á vontade de um Ente supremo. Logo existe \textsc{Deus} \footnote{O consenso do genero humano, relativamente á \emph{existencia} de um Ente supremo, não podia ser mais constante, nem mais universal. Esta crença num Ente supremo podia ser e foi desfigurada por muitos erros; mas conservou sempre o fundo intacto. Muitos povos ignoraram e ignoram ainda a \emph{essencia} de \textsc{DEus}, mas nenhum ignorou ou ignora a sua \emph{existencia}.\par \emph{Cicero} disse: « Não existe povo algum, por inculto e selvagem, que não tenha fé em \textsc{DEus}, embora lhe não conheça a essencia » (\emph{De nat}. \emph{Deor}. I, 17).\par \emph{Quatrefages} escreve: « Obrigado pelo meu ensino a passar revista a todas as raças humanas, procurei o \emph{atheismo} tanto entre os povos mais selvagens, como entre os mais civilisados. Não \emph{o encontrei em parte alguma}, \emph{a não ser nalgum individuo ou eschola muito limitada}, como se viu na Europa no seculo passado e como se vê ainda no presente » (\emph{L'espèce humaine,} pag. 355, Paris, 1877)... E' noutro logar não se exprime menos claramente: « Pouco a pouco fez-se a luz, e os povos da \emph{Australia} e da \emph{Melanesia}, os \emph{Boschimans}, os \emph{Ottentotes}, os \emph{Cafres}, os \emph{Bechuanas}, foram eliminados do numero dos povos atheus e reconhecidos como religiosos » (\emph{Rapport}, pag. 410).\par De resto, os \emph{Annaes da propagação da Fé} provam, com razões claras e baseadas na experiencia, que povos atheus não passam de um mytho.\par Objectam: A crença de todos os povos na existencia de um Ente supremo derivou da \emph{ignorancia}, ou da astucia dos \emph{sacerdotes}, ou da prepotencia dos \emph{legisladores}, ou do \emph{medo}, ou dos \emph{prejuizos}.\par Resposta: A crença unanime de todos os povos em \textsc{Deus} não derivou --- a) nem da \emph{ignorancia}, --- \emph{b}) nem da astucia dos \emph{sacerdotes}, --- \emph{c}) nem da prepotencia dos \emph{legisladores}, --- \emph{d}) nem do \emph{medo}, --- \emph{e}) nem dos \emph{prejuizos}.\par a) Não derivou da \emph{ignorancia}. Por quanto, o estudo das causas naturaes é o meio mais facil e efficaz para se conhecer a existencia de \textsc{DEus}; visto que essas causas, sendo \emph{secundarias}, \emph{contingentes}, \emph{imperfeitas}, exigem uma causa \emph{primeira}, \emph{necessaria}, \emph{perfeitissima}, que é \textsc{Deus}.\par A experiencia confirma a these. Na verdade, os maiores naturalistas são os mais convencidos e enthusiastas adoradores de \textsc{Deus}. --- Citaremos apenas os testemunhos de \emph{Linneu} e \emph{Kepler}.\par \emph{Linneu} escreve: « Quando acordei, \textsc{Deus} \emph{eterno}, \emph{immenso}, \emph{omnisciente}, \emph{omnipotente}, acabava de passar. Vi-o ao longe, e fiquei profundamente maravilhado. Segui os vestigios de seus passos, atravez das obras da creação, e por toda a parte, mesmo nas mais pequenas coisas, que poder! que sabedoria! que ineffavel perfeição! O sol e todo o systema planetario, immenso, incalculavel, appareceram-me suspensos pelo Primeiro Motor, Causa das\pgn{439}causas, Governador e Conservador do universo. Todas as coisas são, pois, testemunho da sabedoria divina; a belleza, a harmonia, as justas proporções das creaturas proclamam o poder do grande \textsc{Deus} » (\emph{Syst}. \emph{de la Nat}. p. 10).\par \emph{Kepler} conclue a sua obra, intitulada « \emph{Os quatro livros da harmonia celeste} », por estas sublimes palavras: « E agora só me resta levantar as mãos e os olhos para ao ceu, e dirigir ao Auctor de toda a luz a minha humilde oração: O' Vós, que, pelos esplendores sublimes que derramastes sobre toda a natureza, elevais os nossos desejos até á divina luz da vossa graça, sede bemdito! Senhor e Creador de todas as consolações, que experimentei no extasis, a que me elevou a contemplação da obra das vossas mãos, eis acabado este livro, que contém o fructo do meu trabalho e de toda a intelligencia, que me destes. Proclamei deante dos homens toda a grandeza das vossas obras, mostrando-lhes, na medida das minhas forças, a sua immensa perfeição. Esforçei-me por chegar até á verdade, por a conhecer o mais perfeitamente possivel, e, se me escapou alguma coisa indigna de Vós, fazei que a conheça para a eliminar. Ter-me-hei deixado seduzir pela presumpção, na presença da belleza maravilhosa das vossas obras? Terei eu procurado a minha gloria entre os homens, levantando este monumento, que só devia ser consagrado á vossa honra? Ah! se assim foi, tende piedade e misericordia de mim, e concedei-me a graça de que o meu trabalho não produza mal algum, mas concorra para a glorificação do vosso Nome e para a salvação das almas ».\par Os mesmos sentimentos encontram-se nos escriptos dos mais celebres naturalistas, como \emph{Copernico}, \emph{Galileu}, \emph{Newton}, \emph{Jussieu}, \emph{Cuvier}, \emph{Agassiz}, \emph{Biot}, \emph{Buffon}, \emph{Lavoisier}, \emph{Brongniart}, \emph{Blainville}, \emph{Bertholet}, \emph{Gay-Lussac}, \emph{Elie de Beaumont}, \emph{Thénard}, \emph{Réaumur}, \emph{Barande}, \emph{Milne-Edwards}, \emph{Dumas}, \emph{Chevreul}, \emph{Liebig}, \emph{Flourens}, \emph{Quatrefages}, \emph{Pasteur}, etc. Nem admira; porque, como diz \emph{Herschel}, « quanto mais se extende o campo da Physica, tanto mais numerosas e\pgn{440}inconcussas se tornam as demonstrações da existencia eterna de uma intelligencia creadora e omnipotente; os geologos, os astronomos, os naturalistas trouxeram todos a sua pedra para este grande templo da sciencia, --- templo levantado ao proprio \textsc{DEUS} ». --- Não nos cançaremos de repetir com \emph{Bacon}; « A pouca sciencia afasta de \textsc{Deus}, a muita conduz a Elle ».\par b) Não derivou da astucia dos \emph{sacerdotes}; pois a instituição dos \emph{sacerdotes é effeito}, e não \emph{causa}, da crença universal na Divindade.\par c) Não derivou da prepotencia dos \emph{legisladores}. Os legisladores basearam as leis na propria crença religiosa dos seus povos. --- Além d'isso, os povos não teriam acceitado, ao menos em toda a parte, uma crença, que é contraria a todas as paixões. --- Finalmente, as leis variam em todos os povos, mas a crença em \textsc{Deus} foi sempre a mesma por toda a parte.\par d) Não derivou do \emph{medo}. Na verdade, esse \emph{medo} ou era conforme com a razão, ou não. Se não era conforme com a razão, não podia produzir a crença em \textsc{DEus}; pois essa crença é universal e constante, e um effeito universal e constante não póde derivar de uma causa particular e mutavel, como é o \emph{medo}. Se era conforme com a razão, esse \emph{medo} já suppunha a crença em \textsc{Deus}; pois ninguem tem \emph{medo} de um ente, que se não conhece como existente. --- A verdade é que o \emph{medo}, longe de produzir a crença em \textsc{Deus}, deu sempre occasião ao \emph{atheismo}, pois os libertinos, \emph{pelo facto de temer o castigo}, esforçam-se por negar a existencia de \textsc{Deus}.\par e) Não derivou dos \emph{prejuizos}; pois os \emph{prejuizos} variam segundo a diversidade dos povos, ao passo que a crença em \textsc{DEus} é universal e constante.}.
